% ============================================================
%  LITERATURE REVIEW
% ============================================================
\section{Background and Related Works}

\subsection{Autoregressive Language Models and the GPT Lineage}

The modern era of large language models can be said to begin with the Transformer architecture
\parencite{vaswani2017attention}, which replaced recurrent networks with self-attention and
positional encodings, enabling vast parallelisation during training.
GPT-2 \parencite{radford2019gpt2} demonstrated that a purely autoregressive, decoder-only
Transformer trained on a large web corpus could achieve strong zero-shot and few-shot performance
across a diverse set of NLP benchmarks without task-specific fine-tuning.
Subsequent scaling--GPT-3 \parencite{brown2020language} and beyond--showed that language
modelling loss and downstream task performance both improve smoothly as a power law in compute, dataset size, and parameter count \parencite{hoffmann2022chinchilla}.
While this scaling paradigm has driven frontier capabilities, it has also substantially raised
the barrier to academic experimentation.

Recent efforts have therefore explored whether carefully designed \emph{small} models can close the gap.
Compact, reproducible frameworks such as NanoGPT \parencite{karpathy2022nanogpt} and its
successor NanoChat \parencite{nanochat} provide minimal, end-to-end implementations of GPT-2
style models that can be trained on a single GPU node for a cost of around \$100. These frameworks enable systematic, reproducible ablation studies that are infeasible at frontier scale.

\subsection{Residual Connections and their Limitations}
While residual connections $h_l = h_{l-1} + f_{l-1}(h_{l-1})$ mitigate vanishing gradients \parencite{he2016deep, vaswani2017attention}, standard residuals accumulate previous outputs with fixed, uniform weights. This causes hidden states to grow $O(L)$ with depth, progressively diluting early-layer contributions---a phenomenon termed \emph{PreNorm dilution} \parencite{chen2026attnres}. To resolve this, \textcite{chen2026attnres} introduced \textbf{Attention Residuals (AttnRes)}, replacing fixed additions with learned, input-dependent softmax attention over preceding layers for selective feature retrieval. We investigate whether these training improvements, proven at large scales, transfer to shallow architectures ($L=12$).

\subsection{Memory-Augmented Language Models}To extend effective context beyond the attention window, \textcite{cheng2026conditional} proposed the \textbf{Engram layer} as a conditional, trainable memory alternative to traditional retrieval methods \parencite{lewis2020rag}. This module augments specific transformer layers with a differentiable key-value store addressed by \(n\)-gram hash IDs. It performs $O(1)$ lookups using the current context, adding retrieved embeddings directly into the residual stream. Initially validated on 27-billion parameter models as a new axis of sparsity, we evaluate whether this memory mechanism yields meaningful performance gains in a significantly data- and parameter-starved regime.

\subsection{Efficient Training: Mixed-Precision and FP8}

Beyond architectural modifications, training precision has emerged as an important efficiency
lever. FP16 mixed-precision training \parencite{micikevicius2018mixed} became standard practice by allowing most operations to run in half-precision while keeping a master copy of weights in FP32 to preserve gradient fidelity.

More recently, the Hopper-generation H100 GPU introduced native FP8 tensor core support, enabling 8-bit floating-point matrix multiplications with theoretical throughput improvements.
\textcite{peng2023fp8} and subsequent work have shown that careful FP8 quantisation during
training can match FP16 performance while approximately doubling arithmetic throughput.
An emerging hypothesis--explored in our experiments--is that the quantisation noise introduced by FP8 may act as an implicit regulariser, improving generalisation in data-limited regimes.

%\subsection{Evaluation of Language Models at Small Scale}

%Evaluating language models reliably at small scales is non-trivial. Standard benchmarks designed for large models often show high variance when applied to models below a certain capability threshold. The CORE evaluation suite from DataComp-LM \parencite{li2025datacomplmsearchgenerationtraining} addresses this by selecting 22 tasks specifically for their low-variance signal at small scales, computing \emph{centered accuracies} to normalise for task difficulty. We adopt CORE as our primary downstream metric, supplemented by validation bits-per-byte (BPB) and tokens per second as efficiency measures.
